\subsection*{Question 1}



Augmenter un nombre de 4~\% revient à le multiplier par \mbox{$1+\dfrac{4}{100}=\np{1.04}$}, \textbf{réponse C}.
















\subsection*{Question 2}



















25\% d'une quantité, c'est le quart de cette quantité. Quatre-vingt c'est quatre fois vingt, donc 25\% de 80 c'est 20, \textbf{réponse B}.
\subsection*{Question 3}



















Dans une heure il y a soixante minutes, un quart d'heure c'est donc 15 minutes, et trois quarts d'heure 45 minutes, \textbf{réponse C}.
\subsection*{Question 4}



















Le premier chiffre différent de $0$ est donc le chiffre des cent-millièmes. Donc ce nombre est \np{0.0000431}, \textbf{réponse A}.
\subsection*{Question 5}



















Puisque \mbox{$5^2=25$}, il faut donc calculer \mbox{$\np{3.14}\times4\times25$}, \textbf{réponse A}.
\subsection*{Question 6}



























































































D'après le graphique 1, la proportion d'hommes parmi les locataires est \np{56.25}\%, \textbf{réponse C}.

\subsection*{Question 7}



















\mbox{$6x-5=4x+3\iff 2x-5=3\iff2x=8\iff x=4$}, \textbf{réponse D}.
\subsection*{Question 8}



















































$f$ est décroissante, on lit que l'ordonnée à l'origine est $2$ (autrement dit \mbox{$f(0)=2$}), et \mbox{$f(3)=0$}, donc la pente de $f$ est \mbox{$\dfrac{f(3)-f(0)}{3-0}=-\dfrac23$}, l'expression de $f$ est donc \mbox{$f(x)=-\dfrac23 x +2$}, \textbf{réponse D}.
